\section{Introduction}
\label{sec:introduction}

% ----------------------------------------------------------------
% TARGET: ~1.5 pages (two-column). Expand each subsection below.
% ----------------------------------------------------------------

\subsection{Motivation}
\label{sec:intro:motivation}

Modern smartphones and wearable devices carry inertial measurement units (IMUs)
composed of a tri-axial accelerometer and a tri-axial gyroscope. These sensors
open the door to applications in pedestrian navigation, gesture recognition,
augmented reality, and biomechanical analysis---all of which require accurate
knowledge of the device's orientation and position in space.

However, each sensor has complementary weaknesses. The accelerometer provides a
long-term stable reference to the gravity vector but is corrupted by any
non-gravitational acceleration and high-frequency vibration. The gyroscope
measures angular velocity with high bandwidth and low noise in the short term,
yet its numerical integration accumulates unbounded drift over time. Neither
sensor alone can deliver a reliable orientation estimate; fusing both is
essential.

% TODO [Fig.~\ref{fig:raw_signals}]: side-by-side plots of raw accelerometer
% and gyroscope signals showing noise vs.\ drift.

\subsection{Problem Statement}
\label{sec:intro:problem}

Given a stream of accelerometer readings
$\mathbf{a}_b(t) \in \mathbb{R}^3$ and gyroscope readings
$\boldsymbol{\omega}_b(t) \in \mathbb{R}^3$, both expressed in the sensor's
\emph{body frame}~$\mathcal{B}$, we seek to compute at every time step:

\begin{enumerate}
  \item A unit quaternion $\mathbf{q}(t)$ representing the orientation of
        $\mathcal{B}$ with respect to a fixed \emph{world frame}~$\mathcal{W}$.
  \item A position vector $\mathbf{p}(t) \in \mathbb{R}^3$ in $\mathcal{W}$,
        obtained by projecting the acceleration into the world frame, removing
        gravity, and double-integrating.
\end{enumerate}

% TODO [Fig.~\ref{fig:frames}]: illustration of body frame attached to the
% sensor vs.\ world frame (N-E-D or E-N-U).

\subsection{Proposed Approach}
\label{sec:intro:approach}

Our approach proceeds in three stages:

\begin{enumerate}
  \item \textbf{Attitude estimation.} A complementary filter blends the
        low-frequency content of the accelerometer (via a low-pass filter) with
        the high-frequency content of the gyroscope (via a high-pass filter) to
        produce a drift-free, low-noise orientation estimate expressed as a unit
        quaternion.
  \item \textbf{Frame rotation.} The quaternion is used to rotate each
        accelerometer sample from the body frame into the world frame using the
        conjugation $\mathbf{a}_w = \mathbf{q} \otimes \mathbf{a}_b \otimes
        \mathbf{q}^*$, after which the constant gravity component is subtracted.
  \item \textbf{Position integration.} The resulting gravity-free
        world-frame acceleration is double-integrated over time to recover
        velocity and position.
\end{enumerate}

% TODO [Fig.~\ref{fig:pipeline}]: block diagram of the full pipeline
% (sensors → filter → quaternion → rotation → integration → position).

\subsection{Document Structure}
\label{sec:intro:structure}

The remainder of this paper is organised as follows.
Section~\ref{sec:reference_frames} defines the coordinate frames and motivates
the need for a rotation formalism.
Section~\ref{sec:quaternions} introduces quaternion algebra and its application
to three-dimensional rotations.
Section~\ref{sec:complementary_filter} derives the complementary filter used
for attitude estimation.
Section~\ref{sec:projection} describes the projection of accelerometer data
into the world frame and the subsequent position integration.
Section~\ref{sec:results} presents experimental results, and
Section~\ref{sec:conclusions} concludes with a summary and directions for future
work.
